\documentclass[11pt,letterpaper]{article}

\usepackage[T1]{fontenc}
\usepackage{amsmath,amssymb}
\usepackage{tikz}
\usepackage{pgfplots}
\pgfplotsset{compat=1.18}
\usepackage{algorithm}
\usepackage{algpseudocode}
\usepackage{booktabs}
\usepackage{enumitem}
\usepackage[hidelinks]{hyperref}

\newcommand{\N}{\mathbb{N}}
\newcommand{\Z}{\mathbb{Z}}
\newcommand{\Q}{\mathbb{Q}}
\newcommand{\R}{\mathbb{R}}
\newcommand{\C}{\mathbb{C}}
\newcommand{\Map}[3]{#1 \colon #2 \to #3}
\newcommand{\Mod}[1]{#1\text{-}\mathrm{Mod}}
\newcommand{\GroupHom}{\mathrm{Hom}_{\mathbf{Grp}}}
\newcommand{\RingHom}{\mathrm{Hom}_{\mathbf{Ring}}}

\newtheorem{theorem}{Theorem}
\newtheorem{lemma}[theorem]{Lemma}
\newtheorem{definition}[theorem]{Definition}
\newtheorem{remark}[theorem]{Remark}
\newenvironment{proof}{\noindent\textit{Proof.}}{\hfill\(\square\)}

\title{Generalized Collatz Library}
\author{ascendnoosphere}
\date{\today}

\begin{document}
\maketitle

\begin{abstract}
Exact exponent gcd rebasing removes the coprimality restriction from a
positive-divisor common-base escape theorem. The new presentation has identical
integer coefficients and preserves every raw branch and seed. Its complete
phase sweep has fewer phases when the exponent gcd exceeds one. We obtain
exact signed rays, explicit bounds for every raw time, eventual agreement with
the offset's sign, and zero as the only periodic point. Composite bases and
nonunit residuals are allowed.
\end{abstract}

\section{Raw dynamics and parameters}

For integers \(Q,C,D\), with \(D\ge2\), define
\[
F(n)=
\begin{cases}
n/D,&D\mid n,\\
Qn+C,&D\nmid n.
\end{cases}
\]
Division has priority, including at zero. Throughout this article,
\begin{equation}\label{Parameters}
\begin{gathered}
Q=t^u,\qquad D=t^v,\qquad C=t^a s,\\
t\ge2,\quad 0<u<v,\quad a\ge u+v,\quad s\in\Z\setminus\{0\}.
\end{gathered}
\end{equation}
The exponents are natural numbers. No coprimality of \(u,v\) is assumed.
The base \(t\) may be composite, and neither the offset residual nor a state
residual is required to be a unit.

The coprime phase theorem is proved in \cite{AscendCommonBase}.
The prime-source forced itinerary in \cite{QcdSigned} uses
\(g=\gcd(u,v)\), \(L=v/g\), and \(E=u/g\). The present result treats its
positive-divisor, positive pure-multiplier intersection, with \(0<u<v\).
It does not assert a negative-divisor transfer or a result for the other
offset regimes or nontrivial residual multipliers.

\section{An identical coprime presentation}

\begin{definition}[Rebased parameters]\label{RebaseDefinition}
Put
\[
\begin{gathered}
g=\gcd(u,v),\qquad T=t^g,\qquad U=u/g,\\
V=v/g,\qquad A=a/g,\qquad S=t^{a\bmod g}s.
\end{gathered}
\]
All quotients and remainders here are natural-number quotients and remainders.
\end{definition}

\begin{lemma}[Admissibility and reconstruction]\label{Reconstruction}
Under \eqref{Parameters},
\[
g>0,\quad T\ge2,\quad 0<U<V,\quad \gcd(U,V)=1,\quad
A\ge U+V,\quad S\ne0,
\]
and
\begin{equation}\label{IdenticalCoefficients}
Q=T^U,\qquad D=T^V,\qquad C=T^A S.
\end{equation}
\end{lemma}

\begin{proof}
Since \(u>0\), its gcd with \(v\) is positive. The gcd divides both exponents,
so \(u=gU\) and \(v=gV\). Since \(g\le u\), \(U\ge1\).
Multiplication by positive \(g\) preserves strict inequalities; hence
\(u<v\) implies \(U<V\). The integer identity
\[
\gcd(gU,gV)=g\gcd(U,V)
\]
therefore gives \(g=g\gcd(U,V)\), and cancellation of \(g>0\) gives
\(\gcd(U,V)=1\). Since \(g\ge1\), \(t^g\ge t\ge2\).

The sum \(u+v=g(U+V)\) is divisible by \(g\). Dividing \(a\ge u+v\) by
positive \(g\), and using monotonicity of the floor quotient, gives
\(A\ge U+V\). The power identities
\[
(t^g)^U=t^{gU}=t^u,\qquad (t^g)^V=t^{gV}=t^v
\]
prove the first two reconstruction equations. Euclidean division gives
\(a=gA+(a\bmod g)\), so
\[
T^A S=t^{gA}t^{a\bmod g}s=t^a s.
\]
Finally, \(t^{a\bmod g}\) is positive and \(s\ne0\), so \(S\ne0\).
Every equality holds in the integers, including negative \(s\).
\end{proof}

\begin{lemma}[Phase and parameter bounds]\label{Bounds}
The rebased parameters satisfy
\[
T\le Q,\qquad |S|\le |C|.
\]
If \(g>1\), then \(V<v\).
\end{lemma}

\begin{proof}
The positive gcd satisfies \(g\le u\). Powers of a base at least two are
nondecreasing, so \(T=t^g\le t^u=Q\).
Since \(a\bmod g\le a\), positivity of the powers gives
\[
|S|=t^{a\bmod g}|s|\le t^a|s|=|C|.
\]
Also \(v=gV\) and \(V\ge1\). If \(g>1\), then \(gV>V\), hence \(v>V\).
When \(g=1\), the instance is already coprime; no strict phase decrease is
claimed.
\end{proof}

Equations \eqref{IdenticalCoefficients} identify the raw functions, not merely
a quotient or a selected sequence of times. For every integer, the divisor,
division test, quotient, and affine expression are identical. The seed is
unchanged. Thus the coprime theorem applies without a state conversion or a
new orbit estimate. There is only one replacement instance, and the bounds
in Lemma~\ref{Bounds} prevent hidden growth of the base and offset residual.

\section{Every rebased phase and its raw clock}

\begin{definition}[Whole-base order]\label{Order}
A nonzero integer \(x\) has whole-base \(T\) order \(k\) if
\(x=T^k w\) for an integer \(w\) such that \(T\nmid w\).
This is not a prime valuation. In particular, \(w\) may share prime factors
with \(T\).
\end{definition}

For \(k<V\), a state of order \(k\) is division-free because \(D=T^V\).
Put
\[
G_V(T)=\sum_{i=0}^{V-1}T^i,\qquad
P=U+V,\qquad
\beta_k=S T^{A-U-V+1+k}G_V(T).
\]

\begin{theorem}[Exact raw ray]\label{RawRay}
Let \(x\) have whole-base \(T\) order \(k<V\).
For every natural \(j\),
\[
F^{jP}(x)=x+j\beta_k,
\]
and the whole-base \(T\) order of this state remains \(k\).
A sweep has \(V\) affine edges and \(U\) division edges.
Moreover
\[
P=\frac{u+v}{g}.
\]
\end{theorem}

\begin{proof}
Lemma~\ref{Reconstruction} supplies exactly the hypotheses of the coprime
raw-ray theorem of \cite{AscendCommonBase}, with base \(T\), exponents \(U,V\),
offset exponent \(A\), and signed residual \(S\). Its whole-base order
condition is precisely Definition~\ref{Order}. That theorem proves the complete
phase permutation, legal raw clocks, and phase-preserving ray for every
\(j\); its increment is the displayed \(\beta_k\).
Equation \eqref{IdenticalCoefficients} identifies its raw map with \(F\) at
every step, so both its endpoint identity and its intermediate itinerary
apply to the original coefficients. Nothing requires a unit residual.
The inherited phase sweep contains \(V\) affine and \(U\) division edges.
Finally, \(u+v=g(U+V)\) and \(g>0\), so division gives the last identity.
\end{proof}

The increment is nonzero: all powers and the geometric sum are positive,
and \(S\ne0\). Its order-preserving property is inherited from the full
coprime theorem, rather than inferred from endpoint equality alone.

\section{Every signed seed and every raw time}

\begin{theorem}[Full-sequence escape]\label{Escape}
For every \(n\in\Z\setminus\{0\}\), \(M\in\N\), and \(r\in\N\),
\begin{equation}\label{EscapeBudget}
|n|+(|n|+1+M D^U)P\le r
\quad\Longrightarrow\quad
M<|F^r(n)|.
\end{equation}
In particular, \(|F^r(n)|\) tends to infinity as \(r\) tends to infinity.
Also
\begin{equation}\label{SignBudget}
|n|+(|n|+1)P\le r
\quad\Longrightarrow\quad
sF^r(n)>0.
\end{equation}
Thus the orbit eventually has the offset's sign.
\end{theorem}

\begin{proof}
Apply the coprime full-raw escape theorem with the parameters of
Lemma~\ref{Reconstruction}. Its initial normalization removes only legal
divisions and requires at most \(|n|\) raw steps. Its subsequent argument
first crosses the signed ray into alignment with \(S\), then bounds every
intervening raw step. The inherited quantitative bound is
\[
|n|+(|n|+1+M(T^V)^U)(U+V)\le r
\quad\Longrightarrow\quad M<|F^r(n)|.
\]
By \eqref{IdenticalCoefficients}, \(T^V=D\), and \(P=U+V\).
This proves \eqref{EscapeBudget} without discarding any raw time.
For each fixed \(M\), its left side is a finite natural threshold, so the
inequality holds for all sufficiently large \(r\). This is the asserted
full-sequence limit, not a limit along sweep endpoints alone.

The inherited sign threshold gives \(S F^r(n)>0\) after
\(|n|+(|n|+1)P\). Since \(S=t^{a\bmod g}s\) and the power is positive, this
is equivalent to \(sF^r(n)>0\). The factor \(t^a\) in \(C=t^a s\) is
also positive, so this is agreement with the offset's sign.
\end{proof}

The sign qualification on the inherited intermediate bound is essential.
Neither \eqref{EscapeBudget} nor its proof assigns an aligned bound to an
opposite-sign initial segment. The initial prefix and the crossing time
are already included in the thresholds.

\begin{theorem}[Periodicity]\label{Periodic}
The point zero is fixed and is the only periodic integer point of \(F\).
\end{theorem}

\begin{proof}
Since \(D\mid0\), division priority gives \(F(0)=0\). For nonzero points,
Lemma~\ref{Reconstruction} allows direct application of the zero-only
periodicity theorem in \cite{AscendCommonBase}; its map is identical to \(F\).
Equivalently, a nonzero periodic orbit repeats a finite set of magnitudes,
contradicting the full-sequence limit in Theorem~\ref{Escape}.
\end{proof}

\section{Nonunit controls}

For \((t,u,v,a,s)=(2,2,4,7,1)\), rebasing gives
\[
(T,U,V,A,S)=(4,1,2,3,2),\qquad (Q,C,D)=(4,128,16).
\]
The carried residual is not a unit modulo the new base. The two phases give
\[
\begin{aligned}
1&\longmapsto132\longmapsto656\longmapsto41,\\
2&\longmapsto136\longmapsto672\longmapsto42,\\
4&\longmapsto144\longmapsto9\longmapsto164.
\end{aligned}
\]
For phase zero, \(\beta_0=2\cdot4\cdot(1+4)=40\); for phase one,
\(\beta_1=160\). Each displayed path has physical raw length three.

A composite original base also works. The parameters
\((t,u,v,a,s)=(6,2,4,7,1)\) give \(T=36\), \(S=6\),
\(Q=36\), \(D=1296\), and \(C=279936\). The path
\[
1\longmapsto279972\longmapsto10358928\longmapsto7993
\]
has increment \(6\cdot36\cdot37=7992\). No prime-base identification is used.

\section{Formal and computational statements}

The Lean module \texttt{CommonBaseRebasing.lean} proves admissibility,
coefficient reconstruction, nonzero residual, exact rays, explicit raw
escape, the full-sequence limit, eventual sign, periodicity, phase decrease,
and parameter bounds. The mathematical hypotheses are \eqref{Parameters};
there is no original coprimality or residual-unit hypothesis.

The independent Go replay constructs both coefficient tuples with exact
signed integers and compares them. It strips only legal initial divisions,
checks every raw branch in two sweeps, counts all affine and division edges,
and tests the intermediate inequality only after alignment.
The finite sample has 807744 certificates, including repeated states and maps;
it is not an infinite-orbit proof. Original bases two through nine,
\(v\le8\), noncoprime exponents, three offset exponents, sixteen signed
residuals, ordinary signed seeds, and forced-phase seeds define that sample.
Zero, invalid parameter contracts, and oversized seeds are refused, not
classified by the sample. Larger exponents and signed initial division
prefixes are separate tests. Cached powers and disjoint reserved buffers have
zero warmed allocations on four measured configurations; no arbitrary-size
performance claim is made.

\bibliographystyle{amsalpha}
\bibliography{CommonBaseRebasing}
\end{document}
